\chapter{Mappare BSPL su MCP}\label{chap:3}

Il confronto condotto nel capitolo precedente ha evidenziato la complementarità tra BSPL ed MCP. Se BSPL modella la struttura logica e causale delle interazioni tra ruoli paritetici astraendo dai dettagli implementativi, MCP fornisce un'infrastruttura concreta a livello applicativo per la trasmissione dei messaggi. Lo scopo del presente capitolo è definire come eseguire protocolli specificati in BSPL impiegando MCP come canale di trasporto, analizzando le problematiche ingegneristiche poste dalla conciliazione tra i due approcci.

La trattazione si articola attraverso l'analisi progressiva di tre modelli architetturali. Il primo modello esplora una mappatura diretta in cui ciascun ruolo del protocollo viene associato esclusivamente a un client MCP o a un server MCP, affidando la trasmissione dei messaggi da server a client all'impiego delle notifiche. Il secondo modello esamina un'architettura centralizzata in cui tutti i partecipanti operano come client MCP e convergono verso un unico server con funzione di snodo per l'instradamento dei messaggi. Il terzo modello, infine, definisce una soluzione decentralizzata basata su nodi ibridi, in cui ogni partecipante include sia un modulo server che un modulo host con uno o più client MCP, consentendo lo scambio diretto dei messaggi e demandando la gestione dello stato locale e dei vincoli causali a un adattatore di protocollo.

\section{Mappatura con ruoli client-server disgiunti}\label{mappatura-disgiunta}

Il primo modello di mappatura consiste nel rendere ciascun ruolo del protocollo BSPL esclusivamente come client oppure come server MCP. In questo schema, lo scambio di messaggi sfrutta due meccanismi diversi a seconda che il ruolo mittente sia un client o un server. Da un lato, i ruoli resi come server espongono un tool per ciascuno schema di messaggio di cui sono destinatari, definendone l'input schema in base ai parametri del messaggio. Quando un ruolo associato a un client deve inviare un messaggio a un server, procede invocando il tool corrispondente tramite il metodo \path{tools/call}. Dall'altro lato, i server MCP non possono avviare autonomamente delle richieste verso i client. Per consentire a un server di inviare un messaggio a un client ci si basa quindi sulle notifiche relative alla modifica del contenuto di una risorsa e sulla successiva lettura della risorsa stessa da parte del client.

\subsection{Criticità del modello}\label{disgiunta-valutazione}

Il meccanismo per la trasmissione dei messaggi impiegato nel modello di mappatura considerato in questa sezione presenta delle gravi limitazioni. Innanzitutto, le notifiche native definite da MCP, come quelle per l'aggiornamento delle risorse (\path{notifications/resources/updated}), non sono pensate per trasportare i parametri di un messaggio. Queste notifiche si limitano a segnalare la modifica dell'URI di una risorsa, senza includere dati nel payload. Di conseguenza, per ottenere i valori dei parametri, il client deve compiere una seconda operazione, effettuando una richiesta di lettura della risorsa tramite \path{resources/read}. Questo meccanismo trasforma l'invio di un singolo messaggio in uno scambio a due passaggi, introducendo latenza e complicando l'interazione. D'altra parte, definire metodi di notifica personalizzati per trasportare direttamente i parametri del messaggio costituirebbe una forzatura dello standard MCP, compromettendo l'interoperabilità con implementazioni standard del protocollo.

Inoltre, l'uso delle notifiche subordina l'invio dei messaggi a un'azione preliminare del destinatario non prevista dal protocollo BSPL. Per poter ricevere notifiche relative a una risorsa esposta da un server, il client deve prima inviare una richiesta con il metodo \path{subscriptions/listen}. In BSPL, invece, un ruolo è abilitato a inviare un messaggio non appena possiede localmente i parametri richiesti in ingresso, senza dover attendere che il destinatario manifesti prima il proprio interesse a riceverli. In questo modo, l'emissione del messaggio da parte del server viene bloccata fino a quando il client non decide di inviare la richiesta di sottoscrizione, introducendo un vincolo temporale che non trova alcuna giustificazione nei requisiti informativi di BSPL.

A queste problematiche si aggiunge la mancanza di garanzie di correttezza sui dati ricevuti dal client. Mentre nelle chiamate ai tool il server verifica la conformità dei parametri rispetto all'input schema dichiarato, quando il client legge i parametri di un messaggio tramite \path{resources/read} il protocollo restituisce il contenuto della risorsa senza effettuare alcuna validazione formale della sua struttura. Inoltre, le notifiche e le letture di risorse non forniscono un canale per segnalare errori al mittente. Se i dati contenuti nella risorsa risultano non validi o in conflitto con lo stato locale del client, come nel caso di un tentativo di sovrascrivere un parametro già vincolato e immutabile, il client non dispone di alcun meccanismo per respingere il messaggio o notificare l'errore al server.

Infine, tale asimmetria compromette la modularità dell'architettura dell'agente prevista dal modello \emph{Local State Transfer}. In tale architettura, l'adattatore di protocollo costituisce un componente uniforme che gestisce le comunicazioni in modo omogeneo per qualunque ruolo. Adottare il modello di mappatura che si considera in questa sezione costringe invece a sviluppare due tipi di adattatore con logiche di comunicazione diverse: un adattatore per i ruoli client, che trasmette tramite chiamate ai tool e riceve tramite sottoscrizioni e letture di risorse, e un adattatore per i ruoli server, che riceve gestendo invocazioni remote ed emette inviando notifiche e aggiornando risorse.

Oltre ai limiti del meccanismo delle notifiche, questo modello di mappatura presenta un'incompatibilità strutturale con i protocolli composti da tre o più ruoli. In MCP la comunicazione può avvenire esclusivamente tra un client e un server. Due ruoli resi entrambi come server non possono scambiarsi dati direttamente, così come non lo possono fare due ruoli resi entrambi come client. Di conseguenza, se in un protocollo due ruoli $A$ e $B$ comunicano reciprocamente, uno di essi deve necessariamente operare come client e l'altro come server. Nel momento in cui compare un terzo ruolo $C$ tenuto a comunicare sia con $A$ che con $B$, però, il modello di mappatura fallisce. Per poter scambiare messaggi con il ruolo client, infatti, $C$ dovrebbe assumere la qualifica di server, mentre per interagire con il ruolo server dovrebbe operare come client. Non potendo assumere simultaneamente entrambe le qualifiche, $C$ può comunicare con uno solo dei due partecipanti.

Questa limitazione impedisce la realizzazione di protocolli multi-ruolo, come la \zcref{lst:bspl_purchasewithdelivery} \emph{PurchaseWithDelivery}, presentata nel capitolo precedente. In questo protocollo, innanzitutto i ruoli \textsc{Buyer} e \textsc{Seller} comunicano tra loro per concordare il prezzo, in seguito \textsc{Seller} invia l'ordine di spedizione \lstinline|ship| al ruolo \textsc{Shipper}, e quest'ultimo trasmette la notifica di consegna \lstinline|deliver| a \textsc{Buyer}. Lo scambio di messaggi forma quindi una catena circolare tra tutti i tre ruoli coinvolti, e qualsiasi suddivisione dei ruoli tra client e server rende impossibile completare l'esecuzione: se \textsc{Buyer} e \textsc{Shipper} sono entrambi client non possono scambiarsi il messaggio di consegna, mentre se \textsc{Seller} e \textsc{Shipper} sono entrambi server non è possibile recapitare l'avviso di spedizione.

Il modello a ruoli client-server disgiunti si rivela pertanto impraticabile. Da un lato, l'uso combinato di notifiche e lettura delle risorse introduce ritardi, assenza di schemi formali e vincoli procedurali ingiustificati nello scambio. Dall'altro, il vincolo che impone a ciascun partecipante di essere solo client o solo server preclude la gestione di protocolli a tre o più ruoli. Per superare l'impossibilità di stabilire connessioni dirette tra ruoli o entrambi client o entrambi server, si esplora nella prossima sezione un modello di mappatura che consiste nel rendere tutti i ruoli come client MCP, introducendo un server centrale incaricato di raccogliere e instradare i messaggi tra i diversi partecipanti.

\section{Mappatura centralizzata con server hub}\label{mappatura-hub}

Il secondo modello di mappatura affronta i limiti di connettività riscontrati nella sezione precedente adottando un'architettura a stella. In questo schema, tutti i ruoli previsti dal protocollo BSPL vengono configurati come client MCP, mentre la gestione della comunicazione viene affidata a un unico server MCP centrale che opera come hub di instradamento. Ciascun client stabilisce una connessione dedicata verso il server centrale, mentre tra i diversi client non esiste alcun canale di comunicazione diretto. L'intero flusso di messaggi tra i ruoli converge quindi sull'hub, che assume il compito di ricevere le comunicazioni dai mittenti e recapitarle ai rispettivi destinatari.

Lo scambio di messaggi tra i partecipanti si articola in due passaggi distinti, separando la tratta dal mittente all'hub da quella dall'hub al destinatario. Per consentire l'invio dei messaggi, il server centrale espone un tool dedicato per ciascuno schema di messaggio nel protocollo BSPL, specificandone i parametri all'interno del relativo input schema.
Quando un ruolo intende emettere un messaggio, il suo client invoca il tool corrispondente esposto dall'hub tramite il metodo \path{tools/call}, passando i parametri del messaggio come argomenti della chiamata. Per quanto riguarda la consegna verso il ruolo destinatario, poiché il server hub non può avviare chiamate dirette verso i client, la trasmissione fa ricorso alla primitiva delle risorse e alle notifiche. In particolare, per ciascun ruolo il server mantiene una risorsa che funge da coda dei messaggi in arrivo. All'atto della ricezione di una chiamata a un tool, l'hub deposita il messaggio nella risorsa del destinatario ed emette una notifica di modifica tramite il metodo \path{notifications/resources/updated}. Il client destinatario, ricevuta la notifica, procede a recuperare il payload del messaggio effettuando una richiesta di lettura tramite \path{resources/read}.

Oltre a fungere da canale di instradamento, il server centrale integra al proprio interno l'adattatore di protocollo di BSPL, assumendo la responsabilità di mantenere e convalidare lo stato dell'interazione. Poiché in BSPL la conoscenza non è condivisa globalmente ma appartiene ai singoli partecipanti, per poter verificare la correttezza delle comunicazioni l'hub deve tracciare lo stato informativo di ciascun ruolo in modo indipendente. 
A tale scopo, per ciascuno schema di messaggio definito nel protocollo il server centrale mantiene due relazioni locali distinte: una relazione associata al ruolo mittente e una relazione associata al ruolo destinatario. Quando un client invoca un tool per trasmettere un messaggio, l'hub intercetta la chiamata e applica le regole di viabilità valutandole rispetto allo stato locale del ruolo mittente, verificando che tutti i parametri \adorn{in} risultino già vincolati nelle relazioni locali di quel ruolo per la chiave di sessione considerata e che nessun parametro \adorn{out} o \adorn{nil} sia già presente. Se la verifica ha successo, l'hub memorizza la nuova tupla nella relazione locale del mittente, la deposita nella risorsa del destinatario e restituisce un esito positivo. La tupla viene invece registrata nella relazione locale del destinatario solo nel momento in cui quest'ultimo esegue la lettura della risorsa tramite \path{resources/read}, rendendo i nuovi \emph{binding} disponibili per le successive comunicazioni di quel ruolo.
Al contrario, se la chiamata viola i vincoli causali del protocollo o tenta di sovrascrivere un parametro che è già stato assegnato oppure che deve rimanere indefinito, l'hub restituisce una risposta di errore, impedendo che comunicazioni non conformi si propaghino verso il destinatario.

Un aspetto essenziale per l'esecuzione di qualunque protocollo riguarda, poi, l'associazione tra i ruoli logici definiti nella specifica e le entità concrete che comunicano a livello implementativo.
In BSPL le interazioni sono modellate a livello astratto tra ruoli, mentre sul piano dell'implementazione, nel modello di mappatura che si sta attualmente considerando, operano dei processi client MCP. L'associazione di tali processi client con i ruoli del protocollo risulta indispensabile per garantire l'integrità dell'interazione, assicurando che ciascun messaggio sia emesso esclusivamente dal partecipante autorizzato a farlo, e per consentire all'hub di instradare correttamente le comunicazioni verso i destinatari opportuni. L'adozione di un'architettura a stella semplifica la gestione della scoperta della rete, poiché i client interagiscono unicamente con il server centrale e non necessitano di individuare gli indirizzi degli altri partecipanti. L'onere del controllo dei ruoli ricade quindi interamente sull'hub, che gestisce una tabella di associazione tra ruoli e chiavi dei client. Per autenticare i client in modo sicuro, dunque, ciascuno di essi include una chiave di autenticazione pre-condivisa all'interno degli header delle richieste HTTP verso il server centrale. A ogni invocazione di un tool, l'hub identifica il client chiamante e confronta il suo ruolo con quello previsto dalla specifica per quel messaggio: se l'identità non corrisponde al ruolo autorizzato, la chiamata viene respinta, impedendo che soggetti non legittimati introducano tuple nell'interazione.

Infine, per quanto riguarda la scelta tra i meccanismi di trasporto standard definiti nella specifica di MCP, l'impiego di \texttt{stdio} risulta inadeguato per la realizzazione dell'hub centrale.
La comunicazione su \texttt{stdio} si basa infatti sui flussi standard di input e di output, che costituiscono un canale di comunicazione punto a punto dedicato a un unico interlocutore. Poiché la specifica di MCP non definisce dei meccanismi per condividere questi canali tra più di due entità, un server che opera su \texttt{stdio} può interagire con un solo client alla volta. Questa limitazione rende impossibile connettere simultaneamente al medesimo hub i diversi client associati ai ruoli del protocollo. Al contrario, il trasporto basato su Streamable HTTP opera come un servizio di rete in ascolto su una porta dedicata. Questa modalità permette a molteplici client di aprire connessioni concorrenti verso il server centrale. In questo modo, l'hub può ricevere le invocazioni dei tool dai diversi mittenti e al contempo gestire le richieste di lettura delle risorse da parte dei destinatari.

La \zcref{fig:architettura_hub} illustra visivamente la struttura del modello di mappatura trattato in questa sezione e la dinamica dello scambio dei messaggi tra i partecipanti.

\begin{figure}[!ht]
    \centering
    \includegraphics[width=0.60\textwidth]{diagrams/architettura_hub.pdf}
    \caption{Rappresentazione schematica dell'architettura di mappatura centralizzata con server hub. Tutti i ruoli operano come client MCP e convergono su un unico server centrale, che gestisce il controllo degli accessi, la validazione tramite l'adattatore BSPL e l'instradamento dei messaggi verso i destinatari.}
    \label{fig:architettura_hub}
\end{figure}

\subsection{Implementazione prototipale e caso di studio}\label{hub-prototipo}

TODO sviluppare la sottosezione

\subsection{Valutazione e limiti del modello}\label{hub-valutazione}

L'architettura basata su un server che funge da hub centrale risolve il problema di connettività che rendeva inapplicabile il modello a ruoli disgiunti. Poiché tutti i partecipanti assumono la qualifica di client MCP e stabiliscono una connessione verso l'hub, il sistema può gestire protocolli composti da un numero qualsiasi di ruoli, inclusi gli scenari con tre o più attori come \emph{PurchaseWithDelivery}. Inoltre, grazie all'integrazione dell'adattatore di protocollo all'interno del server centrale, il modello beneficia di un punto di verifica unico in cui viene validato ogni messaggio, impedendo la propagazione di tuple di parametri non conformi.

Sebbene la presenza dell'adattatore centrale prevenga l'immissione di dati non validi e renda superflua la gestione degli errori in fase di ricezione dei messaggi, il ricorso al meccanismo delle notifiche e delle risorse per la tratta dall'hub al destinatario mantiene alcuni svantaggi già evidenziati nel primo modello (\zcref{disgiunta-valutazione}). La consegna dei messaggi continua a richiedere una procedura a due passaggi, composta dall'emissione da parte del server della notifica di aggiornamento della risorsa e dalla successiva richiesta di lettura da parte del client. Questo doppio scambio introduce latenza incrementando il numero di interazioni sulla rete necessarie per completare ogni singola trasmissione. Permane inoltre l'obbligo per ciascun client di effettuare una sottoscrizione a una risorsa dell'hub, introducendo un vincolo procedurale che non trova corrispondenza nella teoria del BSPL.

A questi aspetti si sommano i limiti propri di qualsiasi architettura centralizzata. L'hub costituisce un singolo punto di fallimento per l'intero sistema, in quanto la sua eventuale indisponibilità interrompe completamente l'esecuzione del protocollo, precludendo qualsiasi comunicazione anche tra partecipanti operativi. Inoltre, la convergenza di tutti i flussi di dati e delle logiche di validazione su un unico nodo crea un collo di bottiglia, che rischia di degradare i tempi di risposta all'aumentare del numero di istanze del protocollo in esecuzione o dei messaggi scambiati. Sul piano concettuale, infine, affidare a un'entità centrale la memorizzazione degli stati locali dei ruoli e l'instradamento delle comunicazioni, equivale di fatto ad introdurre un orchestratore dell'interazione, e ciò contrasta con l'approccio decentralizzato di BSPL, che è progettato per modellare interazioni dirette tra entità autonome e paritetiche.

La mappatura basata su server hub dimostra la possibilità di eseguire qualsiasi specifica BSPL tramite MCP, ma a scapito della decentralizzazione. Per conciliare la piena aderenza ai principi di BSPL con i vincoli dell'infrastruttura di trasporto, la sezione seguente introduce quindi un terzo modello di mappatura a nodi ibridi che elimina la necessità di un orchestratore.

\section{Mappatura decentralizzata con nodi ibridi}\label{mappatura-ibrida}

Il terzo modello di mappatura supera la presenza dell'orchestratore centrale introdotto nel modello precedente, ripristinando la natura distribuita e paritetica delle interazioni in BSPL. Per conciliare la simmetria tra i ruoli prevista dal BSPL con l'asimmetria del modello client-server di MCP, ciascun partecipante viene reso come un nodo ibrido, ossia un processo che integra al contempo un server MCP e un host MCP. Il server rimane costantemente in ascolto per ricevere le comunicazioni in ingresso provenienti dagli altri partecipanti, mentre i client MCP vengono impiegati per trasmettere i messaggi verso i server degli interlocutori.

A differenza dei modelli di mappatura precedenti, questo approccio implementa lo scambio dei messaggi sfruttando unicamente la primitiva dei tool, senza ricorrere alle notifiche e alle risorse. Nello specifico, il server MCP di ciascun partecipante espone uno strumento dedicato per ogni schema di messaggio di cui il proprio ruolo è destinatario nel protocollo. Quando un partecipante deve trasmettere un messaggio, il suo client MCP invoca il tool corrispondente esposto dal server del destinatario, trasmettendo i dati all'interno degli argomenti della chiamata.
Ogni tool è definito in modo da rispecchiare lo schema di messaggio a cui è riferito: tutti i parametri che compongono lo schema, siano essi adornati con \adorn{in} oppure con \adorn{out}, vengono dichiarati come proprietà obbligatorie all'interno dell'input schema del tool.

La verifica dei vincoli causali del protocollo e la memorizzazione dello stato locale sono demandate all'adattatore di protocollo, il quale opera come intermediario tra i moduli MCP e la logica di business di ciascun nodo ibrido secondo i principi del modello \emph{LoST}. In fase di emissione, l'adattatore valuta preventivamente la viabilità del messaggio da trasmettere, verificando nelle proprie relazioni locali che, per la chiave di interazione considerata, tutti i parametri \adorn{in} risultino già vincolati e nessun parametro \adorn{out} o \adorn{nil} sia già presente. Se tale condizione è soddisfatta, l'adattatore memorizza la nuova tupla nella relazione locale del mittente e procede con l'inoltro del messaggio invocando il tool opportuno tramite il client MCP. In fase di ricezione, invece, il server MCP intercetta la chiamata al tool e trasferisce il payload all'adattatore. 
Quest'ultimo controlla che, per ciascun parametro, il valore contenuto nel messaggio coincida con un eventuale valore già assegnato nelle relazioni locali del nodo, per la specifica chiave di interazione.
Se la verifica ha esito positivo e la tupla del messaggio non è già presente nella relazione locale del destinatario, l'adattatore la aggiunge a tale relazione e notifica la logica di business. Contestualmente, il server MCP conclude la transazione restituendo una risposta positiva con l'attributo \lstinline|isError| impostato a \lstinline|false|. Se invece il messaggio viola i vincoli del protocollo, la procedura di inserimento dell'adattatore fallisce e il server MCP restituisce una risposta di errore con \lstinline|isError| impostato a \lstinline|true|, notificando al mittente la non conformità del messaggio rispetto allo stato dell'interazione.

L'esecuzione dei protocolli BSPL richiede inoltre che ciascun nodo ibrido sia in grado di determinare l'indirizzo di rete a cui inviare i messaggi destinati a uno specifico ruolo e di verificare l'identità dei partecipanti che invocano i propri tool. In BSPL le interazioni sono modellate unicamente tra ruoli astratti, mentre a livello implementativo le comunicazioni avvengono tra processi distribuiti identificati da indirizzi IP e porte di ascolto. L'associazione tra i ruoli della specifica e gli endpoint concreti è necessaria sia per instradare le chiamate ai tool verso i destinatari corretti, sia per impedire che entità non autorizzate introducano dei messaggi nell'interazione. 
Nella letteratura sui sistemi multi-agente, la scoperta dei partecipanti viene spesso affrontata basandosi sul concetto di \emph{Directory Facilitator}, un registro centralizzato assimilabile a un servizio di ``pagine gialle'' presso cui gli agenti pubblicano i propri descrittori di servizio e ricercano gli indirizzi degli altri partecipanti.
L'utilizzo di un registro centrale reintrodurrebbe tuttavia un punto unico di coordinamento incompatibile con la natura decentralizzata di BSPL. Per superare questa limitazione, nell'ambito dei protocolli basati sull'informazione, Christie et al. hanno introdotto con il sistema Pippi \autocite{christie2022pippi} un meccanismo di associazione dinamica guidato unicamente dai dati. In questa impostazione, l'identità di rete di ciascun partecipante viene trattata come un normale parametro del protocollo, e viene fissata sfruttando l'assioma di immutabilità di BSPL.
Nel modello di mappatura che si delinea in questa sezione, almeno in questa fase iniziale di esplorazione, si opta per una configurazione \emph{out-of-band} delle associazioni ruolo - identità di rete, per preservare la decentralizzazione senza introdurre la complessità di una gestione a tempo di esecuzione. All'avvio del sistema, ogni nodo carica una tabella locale che associa in modo statico ciascun ruolo al relativo indirizzo di rete e a una chiave di autenticazione. Quando un client MCP trasmette un messaggio, include tale chiave all'interno delle intestazioni HTTP della richiesta. Ciò permette all'adattatore del nodo ricevente di validare l'identità del chiamante e di verificare che corrisponda al ruolo autorizzato a trasmettere lo specifico messaggio, prevenendo al contempo rischi legati alla falsificazione degli indirizzi IP.

Infine, per quanto concerne il meccanismo di trasporto, si rileva l'inadeguatezza di \texttt{stdio}, già riscontrata per il modello di mappatura basato su hub centrale (\zcref{mappatura-hub}), anche per l'architettura a nodi ibridi attualmente in esame. La natura strettamente punto a punto dei flussi standard di input e di output impedisce infatti a un server di dialogare con più di un interlocutore alla volta. Sebbene in questa architettura non sia presente un hub centrale, il modulo server di ciascun nodo deve comunque poter ricevere invocazioni di tool da parte di molteplici client. Tale esigenza esclude dunque l'impiego di \texttt{stdio} a favore di Streamable HTTP, il quale consente l'invio concorrente di messaggi da parte di vari nodi verso il server di uno stesso destinatario.

La \zcref{fig:architettura_ibrida} illustra visivamente la struttura del modello di mappatura a nodi ibridi e la dinamica dello scambio dei messaggi tra due partecipanti.

\begin{figure}[!ht]
    \centering
    \includegraphics[width=\textwidth]{diagrams/architettura_ibrida.pdf}
    \caption{Rappresentazione schematica dell'architettura di mappatura decentralizzata con nodi ibridi. Ciascun partecipante adotta una struttura omogenea che integra sia un client che un server MCP, coordinati dall'adattatore BSPL locale basato sul modello LoST.}
    \label{fig:architettura_ibrida}
\end{figure}

\subsection{Implementazione prototipale e caso di studio}\label{ibrida-prototipo}

TODO sviluppare la sottosezione

\subsection{Valutazione e limiti del modello}\label{ibrida-valutazione}

L'architettura a nodi ibridi consente di conciliare la natura paritetica e distribuita di BSPL con l'infrastruttura client-server di MCP, superando i limiti riscontrati nei due modelli di mappatura esplorati nelle sezioni precedenti. Rispetto alla soluzione basata su hub centrale, l'eliminazione del server unico ripristina la decentralizzazione dell'interazione, rimuovendo il punto unico di fallimento e scongiurando la possibilità che si formi un collo di bottiglia. Inoltre, l'impiego esclusivo dei tool per la trasmissione dei messaggi supera definitivamente le criticità legate all'uso combinato di risorse e notifiche (\zcref{disgiunta-valutazione}). Le consegne dei messaggi si compiono infatti in un'unica interazione diretta tra mittente e destinatario, riducendo la latenza ed eliminando l'obbligo di sottoscrizione preventiva a una risorsa. Altresì, viene completamente delegata ad MCP la validazione sintattica dei messaggi, mediante gli input schema dei tool, prima della verifica dell'adattatore sui vincoli informativi.
Sul piano architetturale, infine, ciascun partecipante adotta una struttura omogenea, permettendo all'adattatore di protocollo di operare in modo uniforme per tutti i ruoli.

Nonostante i vantaggi sopra citati, questo terzo modello di mappatura evidenzia un disallineamento concettuale tra la natura asincrona di BSPL e quella bloccante delle richieste in MCP. Nella teoria di BSPL, la trasmissione dei messaggi è un'operazione priva di attese e puramente unidirezionale, e ciò consente agli agenti di procedere nell'esecuzione senza dipendere dalla reattività degli interlocutori. L'adozione dell'invocazione dei tool come metodo per trasmettere i messaggi vincola invece il mittente all'attesa di una risposta da parte del server, introducendo un accoppiamento temporale a livello di trasporto. Tale sincronicità va tuttavia considerata come un compromesso volto a garantire l'affidabilità delle interazioni nei contesti reali. Poiché il server del destinatario emette la risposta non appena l'adattatore locale ha convalidato il contenuto del messaggio, senza attendere l'elaborazione della logica di business dell'agente, tale risposta assume un ruolo puramente di riscontro sincrono su validità causale e sintattica del messaggio. Questo meccanismo consente al mittente di rilevare tempestivamente violazioni di vincoli o errori di trasmissione, compensando la perdita della completa asincronia con maggiori garanzie di correttezza. A tale vincolo temporale si affianca inoltre la rigidità della configurazione \emph{out-of-band}, che richiede la definizione preventiva di tutti i partecipanti e delle relative chiavi, precludendo l'adesione dinamica di nuovi agenti durante l'esecuzione dell'interazione.

In definitiva, pur a fronte dei compromessi appena evidenziati, l'architettura a nodi ibridi si dimostra una soluzione in grado di preservare la natura decentralizzata di BSPL operando al di sopra dell'infrastruttura MCP. 
Il modello analizzato in questa sezione conclude positivamente l'indagine circa la fattibilità architetturale della mappatura di BSPL su MCP, dimostrando che le coreografie guidate dall'informazione possono essere concretamente trasposte sul protocollo MCP senza richiedere modifiche alle sue primitive. 
L'analisi complessiva dei risultati scaturiti da questo percorso, unitamente alla discussione delle prospettive di ricerca volte a superare i limiti identificati, costituisce l'oggetto del prossimo capitolo.